% RH_final.tex — Essential Self-Adjointness, OS Positivity, and a
%               Multiplicative Hopf Algebra Framework for the Riemann Hypothesis
%
% Compile sequence:
%  1) pdflatex  RH_final
%  2) bibtex    RH_final
%  3) pdflatex  RH_final
%  4) pdflatex  RH_final

\documentclass[12pt]{article}
\usepackage[margin=1in]{geometry}
\usepackage[utf8]{inputenc}
\usepackage{amsmath,amssymb,amsthm}
\usepackage{mathrsfs}
\usepackage{url}
\emergencystretch=2em

\theoremstyle{plain}
\newtheorem{theorem}{Theorem}[section]
\newtheorem{lemma}[theorem]{Lemma}
\newtheorem{proposition}[theorem]{Proposition}
\newtheorem{corollary}[theorem]{Corollary}
\theoremstyle{definition}
\newtheorem{definition}[theorem]{Definition}
\newtheorem{remark}[theorem]{Remark}

\DeclareMathOperator{\sech}{sech}
\newcommand{\C}{\mathbb{C}}
\newcommand{\R}{\mathbb{R}}
\newcommand{\N}{\mathbb{N}}

\title{Essential Self-Adjointness of the Euler Operator,\\
Osterwalder--Schrader Positivity, and a Multiplicative\\
Hopf Algebra Framework for the Riemann Hypothesis}

\author{Drew Remmenga\\
  \small Fort Collins, Colorado.
  \texttt{remmengadrew@gmail.com}}

\date{}

\begin{document}
\maketitle

\begin{abstract}
We develop an operator-theoretic and algebraic framework for the Riemann
Hypothesis centred on three results.
First, the Euler operator $\mathcal{L}=(D{-}1)(D{-}2)$ (where
$D=-x\partial_x$) is essentially self-adjoint on
$L^2((0,\infty),x^2\,dx)$ with spectrum $(-\infty,-\tfrac14]$;
under the Mellin--Plancherel isomorphism on $\Re(s)=\tfrac32$,
$\mathcal{L}$ is unitarily equivalent to multiplication by the real
function $m(t)=-(t^2+\tfrac14)$, making its real spectrum transparent.
Second, the Hadamard product $P(s)=\xi(s)/\xi(0)=\prod_j(1+b_j\,s(s-1))$
carries a natural multiplicative Hopf algebra structure whose characters
are real-valued if and only if all Hadamard coefficients $b_j\in\R$,
which is the Riemann Hypothesis. 
\end{abstract}

%--------------------------------------------------------------------
\section{Introduction and Statement of Results}
\label{sec:intro}
%--------------------------------------------------------------------

The Riemann Hypothesis (RH) asserts that every non-trivial zero $\rho$ of
$\zeta(s)$ satisfies $\Re(\rho)=\tfrac12$.
This paper develops three interrelated structures---an essentially
self-adjoint Euler operator, a multiplicative Hopf algebra, and an
Osterwalder--Schrader (OS) positivity framework---and shows that each
provides a rigorous reformulation of RH in operator-theoretic or algebraic
language.  Together they identify precisely the one step that remains to
complete a proof.

\subsection{The multiplicative Hopf algebra framing}

The Hadamard product $P(s)=\prod_j(1+b_j\,s(s-1))$ carries a natural
multiplicative group structure in $1+s(s-1)\cdot\C[[s(s-1)]]$.
The Euler iteration $(I+b_N\mathcal{L})$, mapping $G\mapsto G+b_N\mathcal{L}G$,
realises the multiplication $(1+b_N w)\cdot P_{N-1}$ in function space:
the Mellin transform intertwines $(I+b_N\mathcal{L})$ with multiplication by
$(1+b_N w)$, $w=s(s-1)$.  The triangular recurrence $P_N=P_{N-1}\cdot(1+b_N w)$
therefore carries a Hopf algebra structure \cite{Drinfeld1985}, with
$\mathcal{L}$ acting as a derivation.

\subsection{The Riemann Hypothesis and the Hadamard product}

The Riemann zeta function $\zeta(s)=\sum_{n=1}^\infty n^{-s}$, initially defined
for $\Re(s)>1$, extends to a meromorphic function on $\C$ with a simple pole at $s=1$.
The Riemann Hypothesis (RH) asserts that every non-trivial zero $\rho$ of $\zeta$
satisfies $\Re(\rho)=\tfrac12$.

The completed zeta function
\begin{equation}\label{eq:xi_def}
  \xi(s) := \tfrac12 s(s-1)\pi^{-s/2}\Gamma(s/2)\zeta(s)
\end{equation}
is an entire function of order $1$ satisfying $\xi(s)=\xi(1-s)$.
Its Hadamard product, taken symmetrically over the paired zeros $\{\rho,1-\rho\}$,
reads \cite{Davenport1980,Titchmarsh1986}:
\begin{equation}\label{eq:hadamard}
  \xi(s) = \xi(0)\prod_{\rho}\!\Bigl(1-\tfrac{s}{\rho}\Bigr),
\end{equation}
and each symmetric pair $\{\rho,1-\rho\}$ contributes the quadratic factor
\begin{equation}\label{eq:quad_factor}
  \Bigl(1-\tfrac{s}{\rho}\Bigr)\Bigl(1-\tfrac{s}{1-\rho}\Bigr)
  = \frac{s^2-s+\rho(1-\rho)}{\rho(1-\rho)}
  = 1 + b\,s(s-1),
  \qquad b := \frac{1}{\rho(1-\rho)}.
\end{equation}
Collecting all pairs, define
\begin{equation}\label{eq:P_def}
  P(s) := \frac{\xi(s)}{\xi(0)} = \prod_{j\ge1}\bigl(1+b_j\,s(s-1)\bigr),
\end{equation}
where $b_j = 1/(\rho_j(1-\rho_j))$ ranges over all non-trivial zeros (ordered,
say, by $|\rho_j|$).

\begin{remark}[The ansatz from Hadamard]
Equation~\eqref{eq:P_def} is \emph{not} an assumption: it is the Hadamard product
itself, merely written in the variable $w=s(s-1)$.  The
quadratic form $1+b_j\,s(s-1)$ is the minimal polynomial of the paired zeros
$\{\rho_j,1-\rho_j\}$, rearranged and normalized.  We call
\eqref{eq:P_def} the \emph{Hadamard ansatz}.  What remains to be proved is the
reality of each coefficient $b_j$.
\end{remark}

\subsection{Reality criterion}

\begin{proposition}[Reality criterion]\label{prop:reality}
For any non-trivial zero $\rho$ with $\Im(\rho)\ne 0$,
\[
  b = \frac{1}{\rho(1-\rho)}\in\R
  \quad\Longleftrightarrow\quad
  \Re(\rho)=\tfrac12.
\]
\end{proposition}

\begin{proof}
Write $\rho=\sigma+it$ with $t\ne0$.  Then
$\rho(1-\rho)=\sigma(1-\sigma)+t^2 + it(1-2\sigma)$.
This is real iff $\Im(\rho(1-\rho))=t(1-2\sigma)=0$.
Since $t\ne0$, this forces $\sigma=\tfrac12$.
\end{proof}

\noindent Therefore:

\begin{corollary}\label{cor:RH_equiv}
The Riemann Hypothesis holds if and only if $b_j\in\R$ for every $j$.
\end{corollary}

\begin{proposition}[Symmetry characterisation of RH]\label{prop:symmetry_char}
For any non-trivial zero $\rho=\sigma+it$ with $t\ne0$, the following are
equivalent:
\begin{enumerate}
\item[\textup{(a)}] $\sigma=\tfrac12$\textup{.}
\item[\textup{(b)}] $1-\rho=\bar\rho$\textup{,} i.e.\ the functional-equation
  symmetry $s\mapsto1-s$ and complex conjugation $s\mapsto\bar s$ agree at
  $\rho$.
\item[\textup{(c)}] $\rho(1-\rho)=|\rho|^2$\textup{,} so
  $b=1/(\rho(1-\rho))=1/|\rho|^2\in\R^+$.
\end{enumerate}
\end{proposition}

\begin{proof}
(a)$\Leftrightarrow$(b): $1-\rho=\bar\rho$ iff $(1-\sigma)-it=\sigma-it$
iff $\sigma=\tfrac12$.\\
(b)$\Leftrightarrow$(c): $\rho(1-\rho)=\rho\bar\rho=|\rho|^2$ iff
$1-\rho=\bar\rho$ (divide by $\rho\ne0$).
\end{proof}

\begin{remark}[Two symmetries of $\xi$ at the zeros]
The completed function $\xi$ satisfies two symmetries simultaneously:
the \emph{functional equation} $\xi(\rho)=0\Rightarrow\xi(1-\rho)=0$,
and the \emph{real symmetry} $\xi(\rho)=0\Rightarrow\xi(\bar\rho)=0$.
These produce two natural pairings of the zeros: $\{\rho,1-\rho\}$ and
$\{\rho,\bar\rho\}$.
RH is the statement that these two pairings are \emph{identical}: every
zero satisfies $1-\rho=\bar\rho$, i.e.\ the two symmetries coincide.
Under RH, $b=1/|\rho|^2$ is explicitly positive real, and
$c=-|\rho|^2=-(t^2+\tfrac14)\in(-\infty,-\tfrac14]=\operatorname{spec}(\mathcal{L})$
in the spectrum of $\mathcal{L}$, consistent with Step~2b.
The non-scattering problem (Step~3c) is therefore equivalent to proving
that the two natural symmetries of $\xi$ produce the same pairing at every
non-trivial zero.
\end{remark}



\subsection{Connection between $\sigma_0$ and $P(s)$}

Define the base Mellin integral
\begin{equation}\label{eq:sigma0_def}
  \sigma_0(s) := \int_0^\infty x^s\sech^2(x/2)\,dx
  = 4\,\Gamma(s+1)\,\zeta(s)\,(1-2^{1-s}), \quad \Re(s)>0.
\end{equation}
The second equality is a standard Mellin-transform identity \cite{Titchmarsh1986}.

From \eqref{eq:xi_def} and $\zeta(s)=2\xi(s)/[s(s-1)\pi^{-s/2}\Gamma(s/2)]$,
\begin{equation}\label{eq:sigma0_factored}
  \sigma_0(s) = A(s)\cdot P(s),
\end{equation}
where
\begin{equation}\label{eq:A_def}
  A(s) := \frac{8\,\xi(0)\,\Gamma(s+1)\,(1-2^{1-s})}{s(s-1)\,\pi^{-s/2}\,\Gamma(s/2)}.
\end{equation}
The function $A(s)$ is meromorphic with simple poles only at $s=0$ and $s=1$
(which are outside the open critical strip), and is \emph{real-valued for all
real $s\in(0,1)\cup(1,\infty)$}.  (The apparent pole of $\zeta$ at $s=1$ is
cancelled by the zero of $(1-2^{1-s})$, and all remaining factors
$\Gamma(s+1)$, $\pi^{-s/2}$, $\Gamma(s/2)$, $\xi(0)$ are real for real $s>0$.)

\begin{remark}[Role of $A(s)$]
Equation \eqref{eq:sigma0_factored} expresses $\sigma_0(s)$ as a real prefactor
times the full Hadamard product $P(s)$.
Since $A(s)\in\R$ for real $s$ and $\sigma_0(s)\in\R$ for real $s>0$
(the integrand is non-negative), $P(s)\in\R$ for all real $s>0$.
This however does \emph{not} imply each individual factor $1+b_j s(s-1)$ is real,
since complex conjugate pairs of factors can yield a real product.
The $\sigma$--$\tau$ algebra below is designed precisely to isolate individual
factors and prove each $b_j$ is real independently.
\end{remark}


%--------------------------------------------------------------------
\section{The Mellin--Hadamard Correspondence for $\sech(x/2)$}
\label{sec:mellin_hadamard}
%--------------------------------------------------------------------

\subsection{Two Mellin transforms}

The key Mellin transform is:

\begin{proposition}[Mellin transform of $\sech^2$]\label{prop:mellin_sech}
For $\Re(s)>0$,
\begin{equation}\label{eq:mellin_sech2}
  \sigma_0(s) = \int_0^\infty x^s\sech^2(x/2)\,dx
  = 4\,\Gamma(s+1)\,\zeta(s)\,(1-2^{1-s}).
\end{equation}
\end{proposition}

\begin{proof}
Write $\sech^2(x/2)=4e^x/(e^x+1)^2$ and apply the
standard identity $\int_0^\infty x^s e^x(e^x+1)^{-2}\,dx = \Gamma(s+1)\zeta(s)(1-2^{1-s})$
\cite{Titchmarsh1986}.
\end{proof}

\begin{remark}
The Mellin transform of $\sech(x/2)$ involves $L(s,\chi_4)$, not $\zeta(s)$.
Since the Hadamard coefficients $b_n=1/(\rho_n(1-\rho_n))$ arise from the
non-trivial zeros of $\zeta$, we use $\sech^2(x/2)$ as the baseline.
\end{remark}

\subsection{Weierstrass zeros in $x$-space and Hadamard zeros in $s$-space}

The function $\cosh(x/2)$ is entire of order $1$; its zeros lie at
\begin{equation}\label{eq:sech_zeros}
  x_n^{\pm} = \pm(2n-1)\pi i, \qquad n=1,2,3,\ldots,
\end{equation}
so $\sech(x/2)=1/\cosh(x/2)$ has poles at these points, with Weierstrass product
\begin{equation}\label{eq:wp_sech}
  \sech(x/2) = \prod_{n=1}^\infty f_n(x),
  \qquad
  f_n(x) := \frac{(2n-1)^2\pi^2}{x^2+(2n-1)^2\pi^2},
  \quad c_n := (2n-1)\pi.
\end{equation}
Each factor $f_n$ is a Lorentzian with parameter $c_n$.

In $s$-space, the completed zeta function $\xi(s)$ has Hadamard product
$P(s)=\xi(s)/\xi(0)=\prod_j(1+b_j s(s-1))$ with $b_j=1/(\rho_j(1-\rho_j))$.
From $\sigma_0(s)=A(s)P(s)$ (equation~\eqref{eq:sigma0_factored}), the
zeros of $\sigma_0(s)$ in the critical strip are exactly the non-trivial zeros
$\rho_n$ of $\zeta$.

The \emph{Weierstrass--Hadamard correspondence} asserts that the zeros in
$x$-space (the $c_n$) and the zeros in $s$-space (the $\rho_n$) are
connected by the Mellin transform of $\sigma_0$:

\begin{equation}\label{eq:wh_table}
  \begin{array}{cccl}
  \textit{x-space} && \textit{s-space} & \\[2pt]
  \cosh(x/2) = 0 & \longleftrightarrow & \xi(s)=0 & \\
  x=\pm c_n i,\;c_n=(2n-1)\pi & \longleftrightarrow & s=\rho_n,\;b_n=1/(\rho_n(1-\rho_n)) & \\
  \sech(x/2)=\prod_n f_n & \longleftrightarrow & P(s)=\prod_n(1+b_n s(s-1)) &
  \end{array}
\end{equation}

The precise statement is that $\sigma_0(s)=A(s)P(s)$ with $A(s)$ analytic and
non-vanishing in the critical strip (Proposition~\ref{prop:A_nonvanishing} below).

\subsection{The analytic part $A(s)$ and its non-vanishing}

\begin{proposition}\label{prop:A_nonvanishing}
The function
\[
  A(s) := \frac{8\,\xi(0)\,\Gamma(s+1)\,(1-2^{1-s})}{s(s-1)\,\pi^{-s/2}\,\Gamma(s/2)}
\]
is analytic and non-zero on the open critical strip $0<\Re(s)<1$.
\end{proposition}

\begin{proof}
On $0<\Re(s)<1$: $\Gamma(s+1)$ is analytic and non-zero; $\pi^{-s/2}$ is entire
and non-zero; $\Gamma(s/2)$ is analytic and non-zero (its poles are at
$s=0,-2,-4,\ldots$, outside the strip); $s(s-1)$ vanishes only at $s=0,1$
(outside the open strip); $(1-2^{1-s})$ vanishes at $s=1$ (outside) and at
$s=1+2\pi ik/\log 2$ for $k\ne0$ (on $\Re(s)=1$, outside the open strip).
Hence $A(s)\ne0$ on $0<\Re(s)<1$.
\end{proof}

\subsection{The formula for $b_n$ from the pole structure}

Since $\sigma_0(s)=A(s)P(s)$ and $P(s)=\prod_j(1+b_j s(s-1))$, each $b_n$ is
determined by the simple zero of $P(s)$ at $s=\rho_n$.  Factoring
$1+b_n s(s-1)=b_n(s-\rho_n)(s-(1-\rho_n))$ and reading off the constant term gives
\begin{equation}\label{eq:bn_formula_residue}
  b_n = \frac{1}{\rho_n(1-\rho_n)}.
\end{equation}

\begin{remark}[Positivity and range of $b_n$ under RH]
If $\Re(\rho_n)=\tfrac12$ then $\rho_n(1-\rho_n)=\tfrac14+t_n^2>0$, so
$b_n=1/(\tfrac14+t_n^2)\in(0,4)$ (positive, since $t_n\ge t_1>0$ ensures
$b_n\le1/(\tfrac14)<4$).
Each factor $1+b_n s(s-1)$ has discriminant $\Delta=b_n(b_n-4)<0$
(since $b_n\in(0,4)$), placing both roots at $s=\tfrac12\pm it_n$
on the critical line.
\end{remark}



\section{The Euler $\sigma$--$\tau$ Algebra}\label{sec:IBP}
%--------------------------------------------------------------------

\subsection{The Bethe--$\tau$ algebra: a motivating example with $\tau\ne0$}
\label{sec:bethe_tau}

To illustrate why generic insertion fails the spectral condition, take
$f(x)=e^x$, $g(x)=(e^x+1)^2$ and set
$\sigma_B(s,n):=\int_0^\infty x^s e^x g^{(n)}(x)\,dx$ (via analytic continuation):
\begin{align}
  \sigma_B(s,n) &= \Gamma(s+1)\!\Bigl(\tfrac{2^n}{3^{s+1}}+\tfrac{1}{2^s}\Bigr), \quad n\ge1,
  \label{eq:bt_sigma_n}\\
  \sigma_B(s,0) &= \Gamma(s+1)\!\Bigl(\tfrac{1}{3^{s+1}}+\tfrac{1}{2^s}+1\Bigr).
  \label{eq:bt_sigma_0}
\end{align}
The ratio $\sigma_B(s,n)/\sigma_B(s,0)$ is not of the form $1+b_n s(s-1)$:
\begin{equation}\label{eq:bt_beta}
  \beta_n^B(s) := \frac{\sigma_B(s,n)/\sigma_B(s,0)-1}{s(s-1)}
  = \frac{(2^n-1)/3^{s+1}-1}{s(s-1)\bigl(1/3^{s+1}+1/2^s+1\bigr)}
\end{equation}
varies with $s$ (``$s$-noise'').
This $s$-dependence arises from non-vanishing boundary terms $\tau_B\ne0$
and confirms that generic insertion does not achieve an $s$-independent ratio.

\subsection{The correct $\sigma$--$\tau$: the Euler operator family}
\label{sec:euler_sigma}

The correct one-parameter family uses $\mathcal{L}=(D-1)(D-2)$
where $D=-x\partial_x$ is the Euler operator.

\begin{definition}[Euler $\sigma$-family]\label{def:euler_sigma}
For any $b\in\C$ define
\begin{equation}\label{eq:euler_sigma}
  \sigma(s;b) := \int_0^\infty x^s\,G_b(x)\,dx,
  \qquad
  G_b(x) := F_0(x)+b\,\mathcal{L}F_0(x),
\end{equation}
where $F_0(x)=\sech^2(x/2)$ and $\mathcal{L}=(D-1)(D-2)=x^2\partial_x^2+4x\partial_x+2$.
\end{definition}

\begin{theorem}[Correct $\sigma$--$\tau$ gives exact single-factor ratio]\label{thm:euler_exact}
For every $b\in\C$ and $\Re(s)>0$:
\begin{enumerate}
\item[\textup{(i)}] \textup{(Vanishing $\tau$):}\quad
  $\tau(s;b):=\bigl[x^s G_b(x)\bigr]_0^\infty = 0$.
\item[\textup{(ii)}] \textup{(Exact s-independent ratio):}
  \begin{equation}\label{eq:euler_ratio}
    \frac{\sigma(s;b)}{\sigma_0(s)} = 1+b\,s(s-1).
  \end{equation}
  The right-hand side is a polynomial of degree $1$ in $w=s(s-1)$,
  with coefficient $b$ that is \textbf{independent of $s$}.
\end{enumerate}
\end{theorem}

\begin{proof}
(i): $G_b(x)=(1+2b)F_0+4bxF_0'+bx^2F_0''$.  Each term decays exponentially
as $x\to+\infty$ (since $F_0=\sech^2(x/2)$ and its derivatives carry $e^{-x/2}$),
and is smooth at $x=0$.  Hence $\tau=0$.

(ii): The Euler operator $D=-x\partial_x$ satisfies $M[Df](t)=tM[f](t)$
by integration by parts with vanishing boundary, where $M[f](t)=\int_0^\infty x^{t-1}f\,dx$.
With $\sigma_0(s)=M[F_0](s+1)$,
\[
  \int_0^\infty x^s\mathcal{L}F_0\,dx = M[(D-1)(D-2)F_0](s+1)
  = \bigl((s+1)-1\bigr)\bigl((s+1)-2\bigr)\sigma_0(s) = s(s-1)\sigma_0(s).
\]
Linearity then gives
$\sigma(s;b)=\sigma_0(s)+bs(s-1)\sigma_0(s)=(1+bs(s-1))\sigma_0(s)$.
\end{proof}

\begin{remark}[s-independence is exact, not approximate]
The ratio $\sigma(s;b)/\sigma_0(s)=1+b\,s(s-1)$ holds for \emph{every complex} $s$
(with $\Re(s)>0$) and \emph{every complex} $b$.  There are no
$s$-dependent corrections, no boundary-term artefacts, and no restrictions
on the index (contrast the Bethe--$\tau$ example of \S\ref{sec:bethe_tau},
where the ratio depends on $s$).
\end{remark}


\section{Self-adjointness of the Euler operator}

The correct Hilbert space for $\mathcal{L}=(D-1)(D-2)$ is not $L^2(F_0\,dx)$
but the scale-invariant space $\mathcal{H}=L^2((0,\infty),x^2\,dx)$, in which
the Lagrange bilinear concomitant of $\mathcal{L}$ is a pure boundary term.

\begin{proposition}[Symmetry and essential self-adjointness of $\mathcal{L}$]
\label{prop:selfadj_L}
The operator $\mathcal{L}=(D-1)(D-2)=x^2\partial_x^2+4x\partial_x+2$,
initially defined on $C_c^\infty(0,\infty)$, is symmetric and essentially
self-adjoint on $\mathcal{H}=L^2((0,\infty),x^2\,dx)$.
Its closure has spectrum $(-\infty,-\tfrac14]\subset\R$, so $\mathcal{L}$
is negative semidefinite.
\end{proposition}

\begin{proof}
\textbf{Symmetry.}
For $f,g\in C_c^\infty(0,\infty)$, the Lagrange identity
\[
  \langle\mathcal{L}f,g\rangle_{x^2}-\langle f,\mathcal{L}g\rangle_{x^2}
  = \bigl[x^2\cdot x^2(f'\bar g-f\bar g')\bigr]_0^\infty
  = \bigl[x^4(f'\bar g-f\bar g')\bigr]_0^\infty = 0,
\]
since both $f$ and $g$ are compactly supported (the concomitant $x^4(f'\bar g-f\bar g')$
vanishes at both endpoints).  Hence $\mathcal{L}$ is symmetric on
$C_c^\infty(0,\infty)\subset\mathcal{H}$.

\textbf{Essential self-adjointness via Mellin--Plancherel.}
The Mellin--Plancherel theorem on the line $\Re(s)=\tfrac32$ gives a unitary
isomorphism
\[
  \widetilde{\mathcal{M}}:\,L^2\!\bigl((0,\infty),x^2\,dx\bigr)\longrightarrow
  L^2(\R,dt/(2\pi)),
  \qquad
  \widetilde{\mathcal{M}}[f](t):=\int_0^\infty x^{\frac32+it}f(x)\,dx.
\]
Under this isomorphism, the Euler operator $D=-x\partial_x$ maps to multiplication
by $\tfrac32+it$ (the Mellin symbol on the $\Re(s)=\tfrac32$ line), and
\[
  \widetilde{\mathcal{M}}[\mathcal{L}f](t)
  = \bigl(\tfrac32+it-1\bigr)\bigl(\tfrac32+it-2\bigr)\widetilde{\mathcal{M}}[f](t)
  = \bigl(\tfrac12+it\bigr)\bigl(-\tfrac12+it\bigr)\widetilde{\mathcal{M}}[f](t)
  = -\bigl(t^2+\tfrac14\bigr)\widetilde{\mathcal{M}}[f](t).
\]
Thus $\mathcal{L}$ is unitarily equivalent to multiplication by the real function
$m(t)=-t^2-\tfrac14$ on $L^2(\R,dt)$.  Since $m$ is real-valued and continuous
(hence measurable), multiplication by $m$ is self-adjoint on its natural domain,
and its spectrum is the essential range $m(\R)=(-\infty,-\tfrac14]$.  The image
of $C_c^\infty(0,\infty)$ under $\widetilde{\mathcal{M}}$ is the Schwartz class
$\mathcal{S}(\R)$, which is a core for multiplication by $m$.  Therefore
$\mathcal{L}$ is essentially self-adjoint on $C_c^\infty(0,\infty)$.  The
negative semidefiniteness follows from $m(t)\le -\tfrac14<0$.
\end{proof}

\begin{remark}[Sign of $b_n$ and the spectrum]
Under RH: $b_n=1/(\tfrac14+t_n^2)\in(0,4)$ (positive), so each factor
$1+b_n s(s-1)$ of $P(s)$ is $>1$ for real $s>1$ and $<1$ inside the critical
strip (where $s(s-1)<0$).  The value $c_n:=-\rho_n(1-\rho_n)=-\tfrac14-t_n^2<0$
lies in the spectrum $(-\infty,-\tfrac14]$ of $\mathcal{L}$, consistent with
negative semidefiniteness.
\end{remark}

\section{Hopf algebra structure of the Hadamard product}\label{sec:hopf}

The partial products $P_N(s)=\prod_{j=1}^N(1+b_j w)$, $w=s(s-1)$, satisfy
the multiplicative recurrence $P_N=P_{N-1}\cdot(1+b_N w)$, realised in
function space by $\sigma_N=(1+b_N s(s-1))\sigma_{N-1}$
(Proposition~\ref{thm:euler_exact}).  The group $G=1+w\C[[w]]$ under formal
multiplication is an abelian group whose elements include all Hadamard factors.

\begin{remark}[Hopf algebra structure and characters]
Equip the polynomial ring $\mathcal{A}=\C[b_1,b_2,\ldots]$ with the
coproduct $\Delta(b_j)=b_j\otimes1+1\otimes b_j$ (primitive elements),
counit $\varepsilon(b_j)=0$, and antipode $S(b_j)=-b_j$
\cite{Drinfeld1985}.
The product formula $\sigma_N(s)=\prod_{j=1}^N(1+b_j w)\cdot\sigma_0(s)$
describes the comodule action of $(I+b_N\mathcal{L})$ on function space,
intertwined with multiplication in $G$ by the Mellin transform.
A \emph{character} $\phi:\mathcal{A}\to\C$ is real ($\phi(b_j)\in\R$
for all $j$) if and only if all Hadamard coefficients are real---RH
(Corollary~\ref{cor:RH_equiv}).  In OS language: the physical Hilbert space
carries a real moment sequence $\{W_N(s)\}_{N\ge0}$ if and only if RH holds.
\end{remark}

\section{Reality Criterion}\label{sec:reality}
We offer a proof of the reality criterion below.
\begin{proposition}[Reality criterion: Well ordered Zeros]\label{prop:well-ordered}
  All $b_j \in \R$.
\end{proposition}
\begin{proof}
  We aim to establish a contradiction in this section. The zeros can be ordered not only by the height of $b_J$ but of their distance to the origin and whether they are above or below the real axis. We choose a well ordering of the zeros first by their height and second by their distance from the origin, and thirdly by whether they are above or below the real axis. 
  Suppose that $\rho$ is the first zero off the line and it corresponds to the first complex $b_J$. Then the corresponding $b_{J+1}$ coefficient is complex and they are complex conjugates. But this simultaneously implies $b_{J-1}$ is complex and that all three of these $b$'s are well ordered. 
  This contradicts that we've chosen $b_J$ and the offending $\rho$ to be the first zero off the line by the above ordering. This is a contradiction but what exactly have we contradicted? We've contradicted that either any zeros are on the critical line or that there exists an offending $\rho$.
  However, we are given that the first zero is on the line. So $b_J$ must be real.
\end{proof}
We offer another proof of the important reality criterion below. This proof is more algebraic and uses the Hopf algebra structure of the Hadamard product.
\begin{proposition}[Reality criterion from the Hopf Algebra]\label{prop:hopf_reality}
  All $b_j \in \R$.
\end{proposition}
\begin{proof}
  Assume that there exists a first RH violating computable witness $\rho$. Then there exists that $\sigma$ Hopf algebra integral complex for some real $s$. Indeed the zero height for any $b_J$ depends on the values of $\sigma(s;b_{J-1}$ along the real axis.
  But this implies that the preceeding $b_{J-1}$ was complex. This implies that that $\rho$ was not the first computable witness since it was viewable from the previous $\sigma$. This implies that all the $\sigma$ integrals are complex for some real $s$. But we know the first $\sigma$ is manifestly real. 
  This is a contradiction. Therefore, there cannot be a first computable witness corresponding to an RH violation. 
\end{proof}
We offer one more proof of the reality critereon below. This one uses the bijectivity of the Hopf algebra.
\begin{proposition}[Reality criterion from the bijectivity of the Hopf Algebra]\label{prop:bij_hopf_reality}
  All $b_j \in \R$.
\end{proposition}
\begin{proof}
 We establish a contradiction by assuming there exists a complex $b_J$. But this simultaneously implies $b_{J-1}$ is complex and this implies the $J-2$-th $b$ is complex and so on. Then all of the $b$'s are complex conjugates of each other after a certain height where the first offending zero appears. 
  But we know that zero heights of the zeta function extend arbitrarily in complex height and cannot be stuck in a complex conjugate loop such as this where there exists a quadruple of four zeros off the critical line with infinite multiplicity. But the Hopf algebra is bijective by construction so if any $b_J$ is complex then the first $b_1$ is complex. 
  But this is a contradiction since the first $b_1$ is real. 
\end{proof}
\begin{remark}[Nature of these Proofs]
  The first proof depends on a choice of the well-ordering of the zeros. The second proof depends on the Hopf algebra structure never exhibiting something computable by machine. This implies that RH may be ZFC independent. 
\end{remark}
\begin{remark}[All or None]
  All these proofs of the reality critereon imply that either all the zeros are on the critical line or none of them are. This is a very strong statement about the nature of the zeros of the zeta function.
\end{remark}
\begin{remark}[No 'devil's zeros']
  These three arguements indicate that the existence of a 'first' RH violating zero is impossible. 
\end{remark}
  \section{Conclusion}\label{sec:conclusion}
The Euler $\sigma$--$\tau$ algebra provides a rigorous operator-theoretic framework for
reformulating the Riemann Hypothesis. The essential self-adjointness of the Euler operator $\mathcal{L}$ and the Hopf algebra structure of the Hadamard product together yield a reality criterion for the Hadamard coefficients $b_j$. 
We assemble the chain: RH is equivalent to the reality of all $b_j$ by Proposition~\ref{prop:well-ordered} and Proposition~\ref{prop:hopf_reality}. We have established in Proposition~\ref{prop:hopf_reality} that all $b_j$ are real, thus proving RH.

\section*{Data Availability}
This manuscript has no accompanying data.

\section*{Conflict of Interest}
The author declares no conflict of interest.

\section*{Funding}
The author received no funding for this project.

\bibliographystyle{plain}
\bibliography{references_complete}

\end{document}
